\section{研究路线图}
\label{sec:roadmap}

表~\ref{tab:roadmap}给出一个基于ICA六层模型的研究路线图，按短期（1--2年）、中期（2--4年）和长期（4--8年）三个阶段组织。
每个议题明确说明（1）要解决什么问题、（2）用什么方法、（3）期望的结果。

\footnotesize
\renewcommand{\arraystretch}{1.35}
\begin{longtable}{>{\small}p{0.06\textwidth}>{\small}p{0.06\textwidth}>{\small}p{0.13\textwidth}>{\small}p{0.25\textwidth}>{\small}p{0.18\textwidth}>{\small}p{0.14\textwidth}}
\caption{基于ICA的研究议题清单}
\label{tab:roadmap}\\
\toprule
阶段 & ICA层 & 议题 & 方法 & 度量指标 & 相关定律/公理 \\
\midrule
\endfirsthead
\toprule
阶段 & ICA层 & 议题 & 方法 & 度量指标 & 相关定律/公理 \\
\midrule
\endhead
短期 & L2 & KV缓存命中率优化 & 语义驱逐策略、KV量化压缩、前缀共享 & $H$, $S$, TTFT & 定律I \\[3pt]
短期 & L3 & 统一上下文编译器 & 热/温/冷分层管理，版本戳与TTL标注 & $W_{\mathrm{eff}}$, 压缩比 & 定律II \\[3pt]
短期 & L4 & Tool ABI标准化 & Schema语义版本化、Capability标注 & 成功率、安全违规率 & 公理5 \\[3pt]
短期 & L5 & 软件工程Agent内核 & ACI接口统一，子Agent专业化分工 & 任务完成率、$E$ & 定律III \\[3pt]
中期 & L3 & 语义页置换策略 & 学习型Eviction策略、摘要回写检索 & $\beta(L)$曲线 & 定律II \\[3pt]
中期 & L5 & 一致性可控共享记忆 & 事件日志（Event Sourcing）、冲突检测 & 记忆正确率 & 公理4 \\[3pt]
中期 & L2 & 异构模型协同 & 草稿模型（Speculative Decoding）、MoE路由 & 端到端时延、$H$ & 定律I \\[3pt]
中期 & L4 & Agent安全架构 & 沙箱隔离、审计日志、能力安全（Capability） & 审计完整性 & 公理5,6 \\[3pt]
长期 & L1--L6 & 智能ISA与可编程技能层 & 任务DSL、受控输出格式、可复用技能包 & 泛化能力 & 公理2,3 \\[3pt]
长期 & L3--L5 & 状态化个体Agent & 长期记忆、经验抽象与策略学习 & 经验利用率 & 定律II,III \\[3pt]
长期 & L1--L6 & 分布式模型原生计算织构 & 记忆复制、一致性层、负载均衡 & 集群利用率 & 定律I,III \\[3pt]
长期 & L5 & 治理层原语 & 概率/确定性双平面、审计日志、权限声明 & 治理覆盖率 & 公理3,5,6 \\[3pt]
长期 & L1--L2 & 基底无关的智能计算栈 & 非硅基底接口泛化、基线智能评分阈值 & 接口兼容率 & 公理2 \\[3pt]
长期 & L4--L5 & 物理世界接口驱动 & 传感器/执行器驱动抽象、仿真环境接入 & 驱动覆盖率 & 公理2,5 \\[3pt]
中期 & 全栈 & 最弱环节评估框架 & Worst-at-k、多维度能力剖析、合规测试套件 & 最低维度得分 & 木桶原理 \\
\bottomrule
\end{longtable}
\renewcommand{\arraystretch}{1.0}
\normalsize

\subsection{短期议题}

短期议题聚焦于已有工程原型的优化和标准化，目标是让现有系统的性能和可靠性达到生产级别。

\textbf{KV缓存命中率优化（L2）}
当前问题：大模型推理的解码阶段是典型的内存带宽受限问题\cite{gholami2024memorywall}，KV缓存的命中率$H$直接决定了推理加速比$S$。
方法：通过语义驱逐策略（如H2O的重击者识别\cite{h2o2023}）、KV量化压缩（如KVQuant的8倍压缩\cite{hooper2024kvquant}、KIVI的非对称2-bit量化\cite{liu2024kivi}）和前缀共享（如PagedAttention\cite{kwon2023pagedattention}、RadixAttention\cite{sglang2024}）来提升$H$。
期望结果：在典型Agent工作负载下将$H$从当前的0.5--0.7提升至0.85以上，使定律I预测的加速比$S$达到5--10倍。

\textbf{统一上下文编译器（L3）}
当前问题：不同Agent框架各有自己的上下文管理策略（MemGPT的虚拟上下文、Letta的有状态Agent、HiAgent的层次化工作记忆\cite{memgpt2023,letta_stateful_2026,hiagent2024}），缺少统一的"上下文编译器"来决定哪些信息以原文加载、哪些以摘要加载、哪些只保留索引。
方法：借鉴操作系统的页置换和编译优化，设计热/温/冷三层的上下文管理策略，并为每层附带的上下文块加上版本戳（version stamp）和时效标注（TTL）。
期望结果：在保持任务完成率不降的前提下，将有效上下文利用率$W_{\mathrm{eff}}/C$提升至少50\%，验证定律II的预测。

\textbf{Tool ABI标准化（L4）}
当前问题：不同Agent框架的工具调用接口各异，MCP和A2A提供了协议层面的标准化，但工具schema本身仍缺少版本化、向后兼容和capability标注的统一规范\cite{mcp_intro_2026,agentprotocols2025}。
方法：为工具schema引入语义版本号（如MCP的schema versioning）、输入输出类型约束、副作用声明和capability标注，使L4层真正成为可编程的"应用二进制接口"（ABI）。
期望结果：工具调用成功率提升，安全违规率下降，为公理5（最小权限）的落地提供接口基础。

\textbf{软件工程Agent内核（L5）}
当前问题：SWE-bench和SWE-agent表明，当前代码Agent在真实GitHub问题上的解决率仍然有限\cite{swebench2023,sweagent2024}。
Agent-Computer Interface（ACI）的设计缺乏统一标准，不同框架的子Agent分工策略各异。
方法：统一ACI规范，设计专门的子Agent分工策略（如代码搜索Agent、编辑Agent、测试Agent），并通过定律III的$E$（编排效率）来量化评估。
期望结果：在SWE-bench等基准上将任务完成率提升，同时将$E$从当前的约0.3--0.5提升至0.6以上。

\subsection{中期议题}

中期议题聚焦于跨层协同和系统级安全，目标是让模型原生计算系统从"单点优化"走向"全栈协同"。

\textbf{语义页置换策略（L3）}
当前问题：传统操作系统的页置换（如LRU、CLOCK）基于确定性地址访问模式，而智能系统的"页置换"需要基于语义相关性来决定驱逐哪些上下文块。
方法：设计学习型eviction策略，根据上下文块与当前任务的语义相关度来决定保留或驱逐，对被驱逐的块生成摘要以便回写检索。
期望结果：绘制出$\beta(L)$曲线的精确形状，为定律II提供定量参数。

\textbf{一致性可控共享记忆（L5）}
当前问题：多Agent协作时共享记忆的一致性管理缺乏统一框架。
不同Agent可能同时修改同一记忆块，导致冲突和信息丢失。
方法：引入事件日志（event sourcing）和冲突检测机制，设计可调节一致性级别（从最终一致到强一致）的记忆访问协议\cite{raft2014,cap2002}。
期望结果：在多Agent基准上实现记忆正确率不低于单Agent基线，为公理4（虚拟化公理）在记忆层面的落地提供实证。

\textbf{异构模型协同（L2）}
当前问题：不同规模的模型在延迟、成本和能力上各有优劣，但当前系统通常只使用单一模型。
方法：通过草稿模型（speculative decoding\cite{leviathan2023speculative}）和MoE路由\cite{switch2022,mixtral2024}实现异构模型协同，让简单任务用小模型、复杂任务路由到大模型。
期望结果：在保持输出质量的前提下降低端到端时延，同时提升KV缓存命中率$H$。

\textbf{Agent安全架构（L4）}
当前问题：随着Agent获得越来越多的外部行动能力（读写文件、运行命令、联网），安全边界变得模糊。
方法：为Agent运行时内建沙箱、审计和能力安全机制\cite{debenedetti2025camel,ironclaw2025,agenticsecurity2025}，确保每个有副作用的操作都经过显式审批并产生可审计的事件记录。
期望结果：实现审计完整性100\%（所有副作用操作均可追溯），为公理5和公理6的落地提供安全基础设施。

\subsection{长期议题}

长期议题聚焦于架构层面的根本性创新，目标是构建真正的模型原生计算体系结构。

\textbf{智能ISA与可编程技能层（L1--L6）}
当前问题：提示词、工具描述和技能包构成了模型原生计算的"软指令"，但它们缺少形式化语义和稳定性。
方法：发展任务DSL（Domain-Specific Language）、受控输出格式和可复用技能包，建立类似传统ISA的稳定接口层。如第~\ref{sec:paradigm}节所论证，这一接口层的稳定性是实现基底无关性的前提---无论底层是硅基GPU、神经形态芯片还是生物计算基底，只要满足智能ISA契约，就应能运行同一套Agent框架。
期望结果：Agent的泛化能力显著提升，使得同一技能包可以在不同模型和运行时之间无缝迁移。

\textbf{状态化个体Agent（L3--L5）}
当前问题：大多数Agent是无状态的——每次会话从零开始，无法积累和复用经验。
方法：为Agent引入长期记忆和经验抽象机制，使其能从历史执行中学习策略、总结模式并主动优化行为\cite{amem2025,agentmemorysurvey2025}。
期望结果：经验利用率（历史经验对当前任务的贡献比例）显著提升，为定律II和定律III提供新的验证维度。

\textbf{分布式模型原生计算织构（L1--L6）}
当前问题：当多个Agent集群跨节点协作时，记忆复制、状态一致性和负载均衡等分布式系统问题尚未被系统化地解决。
方法：借鉴分布式数据库的记忆复制和一致性层设计\cite{spanner2012,raft2014}，为多节点Agent集群提供统一的资源调度和状态管理。
期望结果：集群利用率随节点数线性增长（或至少亚线性增长），验证定律I和定律III在分布式场景下的适用性。

\textbf{治理层原语（L5）}
当前问题：现有Agent系统的治理机制（权限、审计、合规）大多是事后添加的补丁，而非体系结构的一等公民。
方法：将概率/确定性双平面、审计日志和权限声明设计为体系结构的核心原语\cite{arbiteros2025,ironclaw2025}，确保每一层都有对应的治理机制。
期望结果：治理覆盖率100\%（每一层的每个有副作用的操作都受到治理约束），实现公理3、5、6的全栈落地。

\textbf{基底无关的智能计算栈（L1--L2）}
当前问题：第~\ref{sec:paradigm}节讨论了神经形态芯片、生物计算和存算一体等非硅基底的快速发展，但当前ICA接口契约假设了GPU-centric的执行模型。
方法：研究L1--L2接口契约如何泛化以适应不同基底（事件驱动的脉冲计算、模拟存内计算、生物神经网络的异步响应），定义基线智能评分阈值作为核心准入条件。
期望结果：制定跨基底的接口契约规范，使未来的非硅计算核心可以无缝接入现有Agent框架。

\textbf{物理世界接口驱动（L4--L5）}
当前问题：第~\ref{sec:paradigm}节指出智能操作系统的下一个拐点是从纯代码交互扩展到物理世界（机器人、自动驾驶、农业、太空导航），但当前的L4工具协议（MCP、A2A）面向数字世界设计。
方法：研究ICA的L4--L5接口如何扩展以支持传感器、执行器、仿真环境和遥测数据流，使其遵循与MCP工具服务器相同的抽象模式。
期望结果：定义面向物理世界的工具驱动规范，使"添加机械臂驱动"遵循与"添加MCP工具服务器"相同的工作流。

\textbf{最弱环节评估框架（全栈）}
当前问题：第~\ref{sec:paradigm}节引入的木桶原理（最弱环节定律）指出系统智能受限于最弱的能力维度，但当前的评估实践以平均指标为主。
方法：设计基于worst-at-k和多维度能力剖析的评估套件，覆盖推理、安全、工具使用、多语言和物理世界交互，作为智能ISA的合规测试。
期望结果：建立智能计算核心的合规级评估标准，使基底无关性从架构原则变为可执行契约。

值得强调的是，这些议题并不要求新模型先行。
许多问题完全可以在现有模型之上开展研究——例如上下文编译器、Tool ABI标准化和Agent安全架构都是系统层面的工程问题，与基础模型的能力提升解耦。
这与传统计算机史相似：架构突破往往不只来自晶体管更强，也来自接口和运行时更成熟。
ISA的稳定性让软件生态繁荣，虚拟内存的发明让程序不再受物理内存限制，POSIX的标准化让应用可以跨平台迁移。
模型原生计算系统同样需要这些"系统层面的突破"。

从产业层面看，第~\ref{sec:sociotechnical}节的分析表明，这些系统层面的标准化与解耦正在与产业结构的演变相互促进：
协议标准化（MCP、A2A、Tool ABI）降低了层间耦合，为第~\ref{sec:sociotechnical}节预测的fabless AI产业结构创造了技术条件；反过来，产业的专业化分工又加速了每层的独立优化。
这种技术与产业的协同演进，正是半导体和PC产业四十年的核心叙事。

%% ============================================================
